\section{Geometry, numerical error, and approximation}\label{sec:theory}

At a fixed location $x\in I$, the ratio in \cref{eq:threshold} is
the lower bound on $\gamma$ required to satisfy the sign restriction
at that location. Because $f_{\boldsymbol g}$ depends linearly on
$\boldsymbol g$ and $h_r$ is fixed, this bound is a linear function
$\boldsymbol\beta_r(x)^T\boldsymbol g$. Here
$\boldsymbol\beta_r(x)\in\R^m$ is the vector of coefficient
multipliers, with its entries ordered as those of $\boldsymbol g$.
For positivity, for example, $h_0=1$ and
$\boldsymbol g=(a_1,\ldots,a_m)^T$, so
$\boldsymbol\beta_0(x)=(-\phi_{1,L}(x),\ldots,-\phi_{m,L}(x))^T$.
The continuum threshold is therefore
\[
 \tau_r(\boldsymbol g)
   =\max_{x\in I}\boldsymbol\beta_r(x)^T\boldsymbol g.
\]

The continuity assumptions and positive derivative bound in
\cref{thm:finite} make $\boldsymbol\beta_r$ continuous and bounded
on $I$, so the maximum representation makes $\tau_r$ convex,
positively homogeneous, and globally Lipschitz.
The threshold enters the log posterior through the normal survival
weight, so its derivatives are needed to compute the gradients used
in posterior sampling. Danskin's theorem \citep{danskin1966maxmin}
expresses these derivatives through the coefficient vectors at
maximizing locations. When the maximizing location $x_\star$ is
unique, the threshold has gradient $\boldsymbol\beta_r(x_\star)$.
Several maximizing locations also permit a gradient when their
coefficient vectors agree, whereas different vectors produce kinks.
The supplementary material gives the threshold derivatives and the
resulting log-posterior gradient formulas.
For a joint restriction in \cref{prop:joint},
the vectors are formed under the common coefficient map and the
maximum ranges over pairs $(r,x)\in A\times I$. A change in the
maximizing derivative order can therefore produce a kink as well
as a change in its location.

Let $\boldsymbol B\boldsymbol B^T=\boldsymbol V_g$ with
$\boldsymbol B$ invertible, and write
$\boldsymbol g(\boldsymbol z)=\boldsymbol m_g+\boldsymbol B\boldsymbol z$.
The exact whitened log density, up to a constant, is
\begin{equation}\label{eq:whitened-target}
 -\tfrac12\|\boldsymbol z\|^2+
 \log\Phi^{\mathrm c}\left(\frac{\tau_r(\boldsymbol g(\boldsymbol z))
                         -\mu(\boldsymbol g(\boldsymbol z))}{s}\right).
\end{equation}
The supplementary material establishes that this density is
$1$-strongly log-concave, giving a unique mode and finite moments of
every polynomial order.

The continuum maximum generally requires numerical evaluation, and
an underestimated threshold would admit curves that violate the
intended sign restriction. We therefore use a conservative upper
threshold $\widetilde\tau_r\geq\tau_r$, so that imposing
$\gamma\geq\widetilde\tau_r(\boldsymbol g)$ preserves the restriction.
The cost is that admissible coefficient values between the exact and
upper thresholds are excluded from the posterior.
The upper threshold may not be
convex, so the log-concavity conclusion need not carry over to this
numerical posterior. Its difference from the exact law can instead
be quantified through the probability mass of the excluded
coefficients, as the following proposition shows under exact Gaussian
operations.

\begin{theoremEnd}[combined main,category=error]{proposition}[Posterior error from a conservative continuum threshold]
\CombinedResultCheck{prop:error}{proposition}\label{prop:error}
Under \cref{thm:finite}, let $\widetilde\tau_r$ be measurable and satisfy
$\tau_r\leq\widetilde\tau_r\leq\tau_r+\delta$ almost everywhere
under the ordinary posterior of $\boldsymbol g$, for a fixed $\delta\geq0$.
Use the same upper threshold in marginal weighting and reconstruction.
Let $P_{\boldsymbol y}$ be the ordinary coefficient posterior, and use
$\mathcal C_r$ also for its coefficient preimage. Define
$\widetilde{\mathcal C}_r=\{\gamma\geq\widetilde\tau_r(\boldsymbol g)\}$.
Both conditioned laws exist and, with $\TV(P,Q)=\sup_A|P(A)-Q(A)|$,
\begin{equation}\label{eq:tv}
 \begin{split}
 \TV\bigl(P_{\boldsymbol y}(\cdot\mid\widetilde{\mathcal C}_r),
          P_{\boldsymbol y}(\cdot\mid\mathcal C_r)\bigr)
 &=1-\frac{P_{\boldsymbol y}(\widetilde{\mathcal C}_r)}
                 {P_{\boldsymbol y}(\mathcal C_r)}
 \leq\min\left(1,\frac{\delta}
 {s\sqrt{2\pi}P_{\boldsymbol y}(\mathcal C_r)}\right).
 \end{split}
\end{equation}
Every measurable summary of the reconstructed coefficients inherits
this bound. The identity holds for any inner event of positive probability.
\end{theoremEnd}
\begin{proofEnd}
\CombinedProofCheck{prop:error}
Fix the ordinary Gaussian law, selected order, and almost-everywhere
error condition of \cref{prop:error}, modifying the numerical threshold
on the exceptional null set if necessary. Finite thresholds and $s>0$
give positive conditional tails, so both conditioning events have
positive probability. To compare the resulting laws, put
$A=\mathcal C_r$ and $B=\widetilde{\mathcal C}_r$. If the omitted
mass is zero, both conditioned laws agree. Otherwise the law on $A$
is a mixture of the laws on $B$ and $A\setminus B$, with weights
$P_{\boldsymbol y}(B)/P_{\boldsymbol y}(A)$ and its complement.
The components have disjoint supports, so total variation equals the
second weight, establishing the identity for any positive-probability
inner event.

For the threshold events considered here, conditional
Gaussian integration identifies the omitted mass as
\[
 \E\left[\Phi((\widetilde\tau_r(\boldsymbol g)-\mu(\boldsymbol g))/s)
          -\Phi((\tau_r(\boldsymbol g)-\mu(\boldsymbol g))/s)
          \mid\boldsymbol y\right].
\]
Since the conditional density is at most $1/(s\sqrt{2\pi})$ and the
excluded interval has length at most $\delta$, dividing this mass
bound by the probability of the exact event proves \cref{eq:tv}.
The invertible coefficient map preserves total variation, while
measurable summaries cannot increase it, proving the remaining assertion.
\end{proofEnd}


The conditioning probability in \cref{eq:tv} explains why a rare
shape event is more sensitive to a given absolute threshold error.
The bound requires the threshold
guarantee to hold globally or almost everywhere under the ordinary
posterior and quantifies the effect of threshold arithmetic within
the finite model, leaving basis approximation and finite-chain error
as separate quantities.

\subsection{Convergence to the conditioned underlying GP}

We next consider the error introduced by replacing the underlying GP
with a finite HSGP expansion, and how this error depends on padding
and rank.
Increasing $L$ separates the artificial boundary from the physical
interval $I$, whereas $m/L$ controls the frequency cutoff, so increasing
padding at fixed rank can remove variation at short scales. The two
choices must therefore be considered together and for the derivative
used in inference as well as for the response
\citep{riutort2023practicalhsgp}.

Because the conditioning event depends on the response or its selected
derivative between grid points, convergence of covariance blocks on
a finite grid is insufficient for comparing the conditioned laws.
We instead work on the entire physical interval in
$C^r(I)$ with its Borel sigma-field and norm
\[
 \|f\|_{C^r}=\sum_{j=0}^r\sup_{x\in I}
             \left|\frac{\mathrm{d}^jf}{\mathrm{d}x^j}(x)\right|,
 \qquad \frac{\mathrm{d}^0f}{\mathrm{d}x^0}=f.
\]

\begin{theoremEnd}{assumption}[Common underlying and finite Gaussian models]
\CombinedResultCheck{ass:underlying}{assumption}\label{ass:underlying}
Fix $I\subset[c-R,c+R]$ with $R>0$ and the observation model and
intercept law specified in \cref{sec:finite-model}. Let $F=a_0+W$, where $W$ is
a centered stationary GP on $\R$, independent of $a_0$, with covariance
function $k$ and spectral density $S_k$ as defined in
\cref{sec:finite-model}. Use that same spectrum and intercept in the finite models in
\cref{eq:finite-model}, varying only $m$ and $L$ while keeping the
center, interval, spectrum, intercept law, observation locations,
positive-definite noise covariance, and data identical in all models.
\end{theoremEnd}

\begin{theoremEnd}{assumption}[Spectral regularity for the selected order]
\CombinedResultCheck{ass:spectral}{assumption}\label{ass:spectral}
Fix $r\in\{0,1,2\}$. The spectrum $S_k$ is continuous, even, and
strictly positive. There is a regularity parameter $\zeta\in(0,1]$
such that the two functions
\[
 \omega\longmapsto S_k(\omega),
 \qquad \omega\longmapsto\omega^{2r+2\zeta}S_k(\omega)
\]
are locally absolutely continuous on $[0,\infty)$, and each function
and its derivative belongs to $L^1(0,\infty)$.
The weighted spectrum $\omega^{2r+2\zeta}S_k(\omega)$ tends to zero
as $\omega\downarrow0$, and we use this limiting value at the origin.
The frequency exponent is
$2\zeta$ for positivity, $2+2\zeta$ for monotonicity, and
$4+2\zeta$ for convexity.
\end{theoremEnd}


Under the Fourier representation of the stationary covariance,
differentiating $r$ times with respect to each input multiplies
the spectrum by $\omega^{2r}$. The additional factor $\omega^{2\zeta}$
allows us to bound the variance of the difference between the $r$th
derivative values at $x$ and $x'$ by a constant multiple of
$|x-x'|^{2\zeta}$. With $\zeta$ held fixed in the approximation limit,
this increment bound supplies the additional regularity needed for
continuity and tightness in $C^r(I)$.

Let $\Pi_{m,L,\boldsymbol y}^{\mathcal C_r}$ and
$\Pi_{\boldsymbol y}^{\mathcal C_r}$ denote the finite and underlying
posterior laws obtained by weighting their priors by the common
Gaussian likelihood and the indicator of \cref{eq:shape-event}.
For laws with finite second norm moments, define
\[
 W_2(P,Q)=\inf_{\lambda}
 \left(\int\|f-g\|_{C^r}^2\,\lambda(df,dg)\right)^{1/2},
\]
where $\lambda$ ranges over couplings with marginals $P,Q$.

Convergence in $W_2$ means that the finite-model curve $F_{m,L}$ and
the underlying curve $F$, with their respective posterior distributions
conditional on the data and $\mathcal C_r$, can be coupled so that
\[
 \E\left[
 \sup_{x\in I}
 \left|
 \frac{\mathrm{d}^jF_{m,L}}{\mathrm{d}x^j}(x)
 -
 \frac{\mathrm{d}^jF}{\mathrm{d}x^j}(x)
 \right|^2
 \right]\longrightarrow0,
 \qquad j=0,\ldots,r,
\]
as $L\to\infty$ and $m/L\to\infty$. Under this coupling, the
expected squared maximum discrepancy over $I$ vanishes for the
response and each derivative through order $r$.

\begin{theoremEnd}[combined main,category=approximation]{theorem}[Convergence of shape-constrained HSGP posteriors]
\CombinedResultCheck{thm:approximation}{theorem}\label{thm:approximation}
Fix $r\in\{0,1,2\}$ under \cref{ass:underlying,ass:spectral}, and
consider finite models with $L\geq2R$ and $m\geq1$ for $r=0$,
$m\geq4$ for $r=1$, or $m\geq2$ for $r=2$.
The underlying GP has a $C^r(I)$ version, and the finite and underlying
shape events have positive probability and null boundaries in that space.
Consequently, the posteriors conditional on the data and $\mathcal C_r$
exist, with finite moments of every polynomial order in $\|\cdot\|_{C^r}$.
As $L\to\infty$ and $m/L\to\infty$,
\begin{equation}\label{eq:wasserstein}
 W_2\bigl(\Pi_{m,L,\boldsymbol y}^{\mathcal C_r},
          \Pi_{\boldsymbol y}^{\mathcal C_r}\bigr)\longrightarrow0.
\end{equation}
This convergence also holds for the posterior summaries: means of
$\frac{\mathrm{d}^jf}{\mathrm{d}x^j}$ converge uniformly on $I$ for
every $0\leq j\leq r$, and covariances between
$\frac{\mathrm{d}^{r_1}f}{\mathrm{d}x^{r_1}}(x)$ and $\frac{\mathrm{d}^{r_2}f}{\mathrm{d}x^{r_2}}(x')$
converge uniformly on $I\times I$ for every $0\leq r_1,r_2\leq r$.
Under the conditioned underlying law,
$\min_I\frac{\mathrm{d}^rf}{\mathrm{d}x^r}>0$ almost surely.
\end{theoremEnd}
\begin{proofEnd}
\CombinedProofCheck{thm:approximation}
Fix $r$ and any sequence of models satisfying the conditions and limits
in \cref{thm:approximation}.
Write $F_{m,L}$ for a draw from the finite prior and
$k_{m,L}(x,x')=\sum_{j=1}^mS_k(\omega_j)\phi_{j,L}(x)\phi_{j,L}(x')$
for the covariance of its spectral component.
We first establish weak convergence of the priors in $C^r(I)$ and
uniform norm moments. Positive shape probability and a boundary of
prior probability zero will then allow passage to the conditioned
posteriors, whose moment bounds give Wasserstein convergence and
convergence of the stated summaries.

\paragraph{Covariance and increment bounds.}
Put $b=2r+2\zeta$. The two spectral weights in
\cref{ass:spectral} also control the intermediate orders needed for
covariance derivatives. For each integer $0<d\leq2r$, differentiating
$\omega^dS_k(\omega)$ on $[0,1]$ gives integrable terms because $S_k$
is bounded there and its derivative is integrable. On $[1,\infty)$,
write the same function as $\omega^{d-b}(\omega^bS_k(\omega))$.
The multiplier and its derivative are bounded, so the product and its
derivative are integrable. Thus, for $d=0,\ldots,2r$ and $d=b$,
the weighted spectra are locally absolutely continuous and have
themselves and their derivatives in $L^1(0,\infty)$, with the cases
$d=0$ and $d=b$ supplied directly by the assumption. Their finite
values at zero and integrable derivatives also give boundedness.
The derivative-covariance argument of
\citet[Theorem 2 and its proof]{differentialensembles2026} then gives, for
$0\leq r_1,r_2\leq r$, with $d=r_1+r_2$:
\begin{equation}\label{eq:covariance-bound}
 \sup_{x,x'\in I}\left|
 \frac{\partial^{r_1+r_2}(k_{m,L}(x,x')-k(x-x'))}
      {\partial x^{r_1}\partial(x')^{r_2}}\right|
 \leq\frac{b_d}{L}+\frac1\pi\int_{\pi m/(2L)}^\infty
                       \omega^dS_k(\omega)\,d\omega,
\end{equation}
where $b_d<\infty$ is independent of $m$ and $L$.
The cited theorem states stronger conditions on classical derivatives,
but its proof bounds frequency-cell errors by integrals of absolute
derivatives and the reflected term by consecutive frequency differences.
Both steps require only the absolute continuity, integrability, and
boundedness verified above. In \cref{eq:covariance-bound}, the first
term vanishes as the padded boundary recedes, while the second vanishes
as the frequency cutoff increases. The common intercept adds no
covariance error.

For the frequency spacing $\Delta=\pi/(2L)$, with $\omega_j=j\Delta$,
the same frequency-cell estimate bounds the spectral sums by
\[
 \Delta\sum_{j=1}^m(j\Delta)^dS_k(j\Delta)
 \leq\|\omega^dS_k(\omega)\|_1+
 \Delta\left\|\frac{\mathrm{d}(\omega^dS_k(\omega))}{\mathrm{d}\omega}\right\|_1,
 \qquad d\in\{2j:0\leq j\leq r\}\cup\{b\}.
\]
Here the norms are over $[0,\infty)$. The right-hand side is uniformly
bounded because $L\geq2R$ gives $\Delta\leq\pi/(4R)$.
Taking $d=0,2,\ldots,2r$ therefore bounds the derivative variances at
$x_-$, with the intercept adding its fixed variance at order zero.

To control increments of the highest derivative, we use this estimate
with $d=b=2r+2\zeta$. Each differentiated basis function is a sine or
cosine with amplitude $L^{-1/2}\omega_j^r$, so
$|\sin u-\sin v|\leq2^{1-\zeta}|u-v|^\zeta$ and the corresponding
cosine inequality give
\[
 \left|\frac{\mathrm{d}^r\phi_{j,L}}{\mathrm{d}x^r}(x)-
        \frac{\mathrm{d}^r\phi_{j,L}}{\mathrm{d}x^r}(x')\right|
 \leq\frac{2^{1-\zeta}}{\sqrt L}
       \omega_j^{r+\zeta}|x-x'|^\zeta.
\]
By independence of the Gaussian coefficients, the variance of the
$r$th derivative increment is the sum of the squared basis increments
above weighted by $S_k(\omega_j)$. The intercept cancels in increments when $r=0$ and
contributes no derivative when $r\geq1$. Consequently,
\[
 \Var\left(\frac{\mathrm{d}^rF_{m,L}}{\mathrm{d}x^r}(x)-
           \frac{\mathrm{d}^rF_{m,L}}{\mathrm{d}x^r}(x')\right)
 \leq C_\zeta|x-x'|^{2\zeta}
       \Delta\sum_{j=1}^m\omega_j^{2r+2\zeta}S_k(\omega_j),
\]
where $C_\zeta=2^{3-2\zeta}/\pi$ follows from $1/L=2\Delta/\pi$.
The frequency-cell estimate at $d=b$ bounds the spectral sum on the
right uniformly in $m,L$, giving the increment bound
$A|x-x'|^{2\zeta}$ with $A$ independent of rank and padding.
Fourier inversion gives the same bound for the underlying mean-square
$r$th derivative, after increasing $A$ if needed. Gaussian moments
then verify the Kolmogorov continuity criterion
\citep{kallenberg2021foundations} with any exponent $p$ satisfying
$p\zeta>1$.
For $r=0$, this gives a continuous version of $W$. For $r\geq1$,
integrating the continuous highest derivative, with the lower
mean-square derivative values at $x_-$ as integration constants,
recovers a version of $W$ by the mean-square fundamental theorem of
calculus. Adding the intercept therefore gives a $C^r(I)$ version of $F$.

\paragraph{Tightness and norm moments.}
We next control both the size and the continuity of the sample paths
uniformly over the finite models. These bounds establish tightness
in $C^r(I)$ and provide the norm moments needed later for Wasserstein
convergence and convergence of posterior summaries.
Let $H$ denote either $\frac{\mathrm{d}^rF_{m,L}}{\mathrm{d}x^r}$ or
$\frac{\mathrm{d}^rF}{\mathrm{d}x^r}$, and choose
$p>\max(4,1/\zeta)$ and a H\"older exponent $0<\kappa<\zeta-1/p$.
For $j\geq0$, let $M_j$ be the largest absolute increment of $H$
between neighboring points on the mesh dividing $I$ into $2^j$
equal cells. The Gaussian increment bound gives
$\E[M_j^p]^{1/p}\leq C_p2^{-(\zeta-1/p)j}$ uniformly over the models,
because the $p$th power of the maximum is bounded by the sum of the
$p$th powers of the $2^j$ increments.
Approximating arbitrary $x,x'\in I$ by points on successively finer
meshes expresses $H(x)-H(x')$ through adjacent increments.
This dyadic chaining argument bounds the H\"older seminorm by
\[
 [H]_\kappa
 =\sup_{\substack{x,x'\in I\\x\ne x'}}
      \frac{|H(x)-H(x')|}{|x-x'|^\kappa}
 \leq C_{I,\kappa}\sum_{j=0}^{\infty}2^{\kappa j}M_j.
\]
Here $C_{I,\kappa}$ depends only on $I$ and $\kappa$.
Minkowski's inequality therefore bounds $\E[[H]_\kappa^p]^{1/p}$
by a constant times $\sum_{j\geq0}2^{-(\zeta-1/p-\kappa)j}$.
Our choice of $\kappa$ makes this series converge, giving a uniform
$p$th moment bound for the seminorm.

The seminorm and the value $H(x_-)$ together bound the highest
derivative everywhere on $I$. When $r\geq1$, repeated integration,
starting from the lower derivative values at $x_-$, then bounds
every lower derivative and the curve itself. For $r=0$, the first
bound already controls the curve. The uniformly bounded Gaussian
moments of these endpoint values consequently give
\begin{equation}\label{eq:moments}
 \sup_{m,L}\E[\|F_{m,L}\|_{C^r}^p]+\E[\|F\|_{C^r}^p]<\infty
\end{equation}
over the permitted ranks and padding. Since $p$ can be arbitrarily
large, every polynomial norm moment is uniformly bounded.

Bounds on derivative values of orders $0,\ldots,r$ at $x_-$ and on the highest
derivative's H\"older seminorm define relatively compact subsets of
$C^r(I)$ by Arzel\`a--Ascoli. The moment bounds and Markov's inequality
allow these bounds to be chosen so that the finite priors assign
arbitrarily high probability to such a set, uniformly over the models.
This proves tightness. To identify the limit, the common mean and
\cref{eq:covariance-bound} give convergence of the joint Gaussian laws
of every finite collection of derivative values. Such evaluations on
a countable dense subset of $I$ determine a $C^r(I)$ law, so every
subsequential limit is the underlying prior. The finite priors
therefore converge weakly to that law.

\paragraph{Positive shape probability and a null boundary.}
The boundary of $\mathcal C_r$ in $C^r(I)$ consists exactly of curves
with $\min_I\frac{\mathrm{d}^rf}{\mathrm{d}x^r}=0$. Indeed, this minimum is continuous
in the $C^r$ norm, and adding or subtracting a small multiple of
$(x-c)^r/r!$ moves any such curve to either side of the event.
We therefore need positive probability of nonnegativity and zero
probability of equality. For each order, these follow by separating
one Gaussian coordinate whose contribution to the selected derivative
is positive throughout $I$.

For $r=0$, the independent intercept gives
\[
 \min_{x\in I}F(x)=a_0+\min_{x\in I}W(x).
\]
Conditional on $W$, nonnegativity therefore imposes a finite lower
threshold on the nondegenerate Gaussian $a_0$, giving positive
conditional probability and zero conditional probability of equality.

For $r\in\{1,2\}$, the intercept contributes no derivative. We instead
seek a decomposition of $W$ into a Gaussian coefficient multiplying
a deterministic function whose $r$th derivative is uniformly positive,
together with an independent Gaussian process. This reduces the shape
restriction to a scalar threshold, as in the finite model.
The reproducing kernel Hilbert space (RKHS) $\mathcal H$ of $W$ in
$C^r(I)$ identifies the deterministic directions that can be separated
in this way.

To construct a candidate direction, choose
$0<\omega_-<\omega_+<\pi/(2R)$ and set
\[
 v_r(x)=\int_{\omega_-}^{\omega_+}S_k(\omega)
             \cos(\omega(x-c)-r\pi/2)\,d\omega.
\]
Taking $r$ derivatives removes the phase shift, leaving the cosine
factor $\cos(\omega(x-c))$. Because $|x-c|\leq R$ and
$\omega_+R<\pi/2$, this factor is uniformly positive on the selected
frequency band and throughout $I$. Thus
\[
 \frac{\mathrm{d}^rv_r}{\mathrm{d}x^r}(x)\geq
 \cos(\omega_+R)\int_{\omega_-}^{\omega_+}\omega^rS_k(\omega)\,d\omega>0.
\]

We next verify that this direction belongs to $\mathcal H$, which
justifies separating a Gaussian coefficient along it.
The Fourier expression for the full-line RKHS norm gives squared norm
$\pi\int_{\omega_-}^{\omega_+}S_k(\omega)\,d\omega<\infty$ for $v_r$,
and restricting the function to $I$ cannot increase that norm.
The RKHS obtained from the covariance on $I$ agrees with the RKHS
of $W$ as a $C^r(I)$-valued Gaussian random element by
\citet[Lemma 8.1]{vdv2008rkhs}, applied to the continuous, injective
embedding into $C(I)$. The norm calculation therefore establishes
$v_r\in\mathcal H$ in the function space used in this proof.

Since $v_r$ is nonzero, we can normalize it by its RKHS norm.
The Gaussian RKHS isometry \citep[Lemma 2.3]{vdv2008rkhs} then supplies
a standard normal variable $\xi$, jointly Gaussian with $W$, such that
$\E[\xi W(x)]=v_r(x)/\|v_r\|_{\mathcal H}$ for every $x\in I$.
Define the remaining process $V$ by
\[
 W=\xi\frac{v_r}{\|v_r\|_{\mathcal H}}+V.
\]
Subtracting this component preserves $C^r(I)$ paths, and
$\Var(\xi)=1$ gives $\Cov(\xi,V(x))=0$ for every $x\in I$.
Joint Gaussianity therefore makes $\xi$ independent of the whole
process $V$.

Conditional on $V$, only the scalar coefficient $\xi$ remains random
in the derivative constraint. Dividing each derivative inequality by
the uniformly positive derivative of $v_r$ gives
\[
 \{F\in\mathcal C_r\}=
 \left\{\xi\geq\|v_r\|_{\mathcal H}
       \max_{x\in I}\frac{-\frac{\mathrm{d}^rV}{\mathrm{d}x^r}(x)}
                              {\frac{\mathrm{d}^rv_r}{\mathrm{d}x^r}(x)}\right\}.
\]
For each path of $V$ the threshold is finite, so the conditional normal
tail has positive probability and no atom at equality. Equality
is precisely the zero-minimum event identified above, so the underlying
shape event has positive probability and a null boundary at all three
orders. The finite laws have the same properties by
\cref{cor:positive,cor:monotone,cor:convex}, applied with
likelihood one. In both cases, conditioning on nonnegativity therefore
gives a strictly positive minimum selected derivative almost surely.

\paragraph{Likelihood and shape conditioning.}
Write $\mathcal L_{\boldsymbol y}(f)$ for the Gaussian density at
$\boldsymbol y$ with mean $(f(x_1),\ldots,f(x_n))^T$ and covariance
$\boldsymbol\Sigma$, taking $\mathcal L_{\boldsymbol y}(f)=1$ when
$n=0$. This likelihood is bounded, positive, and continuous on
$C^r(I)$. For any bounded continuous test function
$\psi:C^r(I)\to\R$, the map
\[
 f\longmapsto\psi(f)\mathcal L_{\boldsymbol y}(f)
                    \boldsymbol1_{\mathcal C_r}(f)
\]
is bounded and continuous outside $\partial\mathcal C_r$, where
$\boldsymbol1_{\mathcal C_r}(f)$ is one if $f\in\mathcal C_r$ and
zero otherwise. The preceding argument shows that this boundary
has zero probability under the underlying prior, so the Portmanteau theorem
\citep[Section 2]{billingsley1999convergence}
gives convergence of the integrals of this map under the finite priors
to its integral under the underlying prior.
Taking $\psi\equiv1$ shows that the normalizing constants converge
to a positive limit. Dividing by them gives convergence of posterior
expectations of every bounded continuous $\psi$, which is conditioned
weak convergence. The posterior densities relative to the priors are
uniformly bounded along the sequence by the likelihood supremum
divided by these constants, so
\cref{eq:moments} also holds for the conditioned laws at every fixed
polynomial order. Likelihood weighting preserves the zero boundary
probability, giving the strictness assertion.

\paragraph{Wasserstein convergence and posterior summaries.}
We now strengthen conditioned weak convergence to convergence in
Wasserstein distance. Index the chosen approximation sequence by
$\ell$, with ranks $m_\ell$ and padding $L_\ell$. In this paragraph,
write $F_\ell$ and $F$ for draws from the finite and underlying
conditioned posteriors, respectively. Since $C^r(I)$ is separable
and complete, the Skorokhod representation theorem
\citep[Theorem 6.7]{billingsley1999convergence} allows these draws
to be constructed on a common probability space, preserving their
posterior laws, so that $\|F_\ell-F\|_{C^r}\to0$ almost surely.
The uniform fourth-moment bounds established above make the squared
norm differences uniformly integrable, so their expectations also
converge to zero. Wasserstein distance takes the infimum over all
couplings, and this particular coupling therefore gives
\[
 W_2^2\bigl(\Pi_{m_\ell,L_\ell,\boldsymbol y}^{\mathcal C_r},
            \Pi_{\boldsymbol y}^{\mathcal C_r}\bigr)
 \leq \E[\|F_\ell-F\|_{C^r}^2]\longrightarrow0,
\]
proving \cref{eq:wasserstein}.

The same coupling controls posterior summaries uniformly in their
input locations because the $C^r$ norm bounds every derivative through
order $r$ over the whole interval. For each such order, the uniform
difference of derivative means is at most
$\E[\|F_\ell-F\|_{C^r}]$, which tends to zero by Cauchy--Schwarz.
For any two derivative orders through $r$, expand the difference of
their products by changing one factor at a time. The resulting
difference of mixed second moments is bounded, uniformly on
$I\times I$, by
\[
 \E[\|F_\ell-F\|_{C^r}(\|F_\ell\|_{C^r}+\|F\|_{C^r})].
\]
Cauchy--Schwarz bounds this expectation by the product of
$\E[\|F_\ell-F\|_{C^r}^2]^{1/2}$, which tends to zero, and
$\E[(\|F_\ell\|_{C^r}+\|F\|_{C^r})^2]^{1/2}$, which remains bounded.
Thus the mixed second moments converge uniformly. Subtracting the
uniformly convergent products of means proves covariance convergence.
Since the sequence was arbitrary, all these conclusions hold in the
full stated limiting regime.
\end{proofEnd}


The spectral weight controls the highest derivative in the
function-space norm, so positivity needs only a continuous process,
whereas monotonicity and convexity require $C^1$ and $C^2$ regularity,
respectively. The proof adapts the derivative-covariance argument of
\citet[Theorem 2 and its proof]{differentialensembles2026} to the
weighted-spectral conditions above, then uses increment moments to
obtain convergence in that space and passes through likelihood
and event conditioning. This last step requires positive event
probability and a null boundary, which the independent random intercept
supplies for positivity and a positive reproducing kernel Hilbert space
(RKHS) derivative direction supplies for the two derivative constraints.

Squared exponential spectra satisfy all three sets of conditions. A Mat\'ern
spectrum is proportional to
$(\lambda^2+\omega^2)^{-(\nu+1/2)}$ with $\lambda>0$
\citep[Chapter 4]{rasmussen2006gpml}. Its weight of order $d$ is
integrable when $d<2\nu$, and the weighted derivative is integrable
for the orders used here, so any
$0<\zeta<\min(1,\nu-r)$ verifies \cref{ass:spectral} when
$\nu>r$. The conditions also hold for positive finite mixtures with
a common smaller exponent. A squared exponential covariance multiplied
by $\cos(\omega_0u)$, with fixed $\omega_0\geq0$, also qualifies
because its spectrum is a positive sum of translated Gaussian densities.
